\documentclass[11pt]{article}

\usepackage[a4paper,margin=1in]{geometry}
\usepackage{amsmath,amssymb,amsthm,mathtools}

\newtheorem{theorem}{Theorem}
\newtheorem{lemma}{Lemma}
\newtheorem{remark}{Remark}

\newcommand{\Pcal}{\mathcal P}
\newcommand{\Fcal}{\mathcal F}
\newcommand{\Bcal}{\mathcal B}
\newcommand{\Ccal}{\mathcal C}
\newcommand{\Tcal}{\mathcal T}
\newcommand{\Ncal}{\mathcal N}
\newcommand{\Om}{\Omega}
\newcommand{\Th}{\Theta}
\newcommand{\1}{\mathbf 1}
\newcommand{\abs}[1]{\left|#1\right|}
\newcommand{\norm}[1]{\left\|#1\right\|}

\title{Non-uniform Non-Archimedean Estimate for Regular Domains}
\author{}
\date{}

\begin{document}
\maketitle

\section*{Statement}

Let \((X,\mu,\Pcal)\) be a good grid and let
\[
  0<s<\frac1p,\qquad 1\le p< \infty, 1\le q\leq  infty
\]
For \(Q\in\Pcal^k\), write \(|Q|=\mu(Q)\), and use the Souza atom
\[
  \psi_Q = |Q|^{s-\frac1p}\1_Q .
\]
There is $C_{na}\geq 0$ such that the following holds: Let \(f\in \Bcal^s_{p,q}\cap L^\infty\) be represented 
\[
  f=\sum_{k\ge0}\sum_{Q\in\Pcal^k} c_Q\psi_Q ,
\]
by either  an standard representation or a positive repsentation,
with coefficient finite  cost
\[
  \norm{c}_{b^s_{p,q}}
  =
  \begin{cases}
  \displaystyle
  \left(
    \sum_{k\ge0}
    \left(
      \sum_{Q\in\Pcal^k}\abs{c_Q}^p
    \right)^{q/p}
  \right)^{1/q},
  & q<\infty,\\[2ex]
  \displaystyle
  \sup_{k\ge0}
  \left(
    \sum_{Q\in\Pcal^k}\abs{c_Q}^p
  \right)^{1/p},
  & q=\infty.
  \end{cases}
\]
Assume also that
\[
  \abs{f(x)}\le M
  \qquad\text{for a.e. }x.
\]

Let \((\Om_i)_{i\in\Lambda}\) be pairwise disjoint domains, such that $1_{\Omega_i}$ with a positive representation 
\[
  \1_{\Om_i}
  =
  \sum_{k\ge0}\sum_{Q\in\Pcal^k} b^i_Q\psi_Q
\]
such that
\[
  b^i_Q\neq0 \implies  b^i_Q> 0\implies Q\subset \Om_i,
  \qquad
  \abs{b^i_Q}\le \Ccal_i,
\]
and such that  $\norm{c}_{b^s_{p,q}}  \leq \Ccal_i$

Let \((\Th_i)_{i\in\Lambda}\) be real weights.  The main
hypothesis is that, for every active source cell \(P\), that is, for every
\(P\in\Pcal\) with \(c_P\neq0\),
\begin{equation}\label{eq:overlap}
  \sum_{\substack{i\in\Lambda\\ P\cap\Om_i\neq\emptyset}}
  \abs{\Th_i}\bigl(1+\Ccal_i\bigr)
  \le N .
\end{equation}
Define
\[
  h
  =
  \sum_{i\in\Lambda}\Th_i\1_{\Om_i}f .
\]
Then there is a representation \(S\) of \(h\), supported in the domains
\(\Om_i\), such that
\begin{equation}\label{eq:main-bound}
  \operatorname{pqCost}(S)+\norm{h}_{L^\infty}
  \le
  C_{\mathrm{na}}\,N\,
  \bigl(\operatorname{pqCost}(R)+M\bigr),
\end{equation}
where \(R=(c_Q)_Q\) is the chosen representation of \(f\), and \(C_{\mathrm{na}}\)
depends only on the structural constants of the grid and on \(s,p,q\).
Moreover, if a coefficient of \(S\) at a cell \(Q\) is nonzero, then
\[
  Q\subset \Om_i
  \qquad\text{for some }i\in\Lambda.
\]

\section*{The Product Decomposition}

Suppose first that $\Lambda$ is finite. If 
\[
  \1_{\Om_i}
  =
  \sum_Q b^i_Q\psi_Q
  \qquad\text{and}\qquad
  f
  =
  \sum_P c_P\psi_P.
\]
The product is decomposed into two parts:
\[
 \left(  \sum_{i\in \Lambda} \Theta_i \1_{\Om_i}\right) f=\sum_{i\in \Lambda}  U_{1,i}+\sum_{i\in \Lambda} U_{2,i}.
\]
The first term \(U_{1,i}\) collects the interactions where the indicator cell is
below the source cell.  If \(Q\in\Pcal^j\), its coefficient is
\begin{equation}\label{eq:u1-coeff}
  (U_{1,i})_Q
  =
  \Theta_i b^i_Q\,\Tcal_R(Q),
\end{equation}
where
\begin{equation}\label{eq:tower}
  \Tcal_R(Q) =
  \sum_{\ell\le j}
  \sum_{\substack{P\in\Pcal^\ell\\ Q\subset P}}
  c_P\,|P|^{s-\frac1p}.\end{equation}

Since the representation of $f$ is either standard or positive  we have 
$$\|\Tcal_R(Q)\|\leq M.$$


This is the ancestral tower of the source representation \(R\), evaluated on
the cell \(Q\).  In Lean this is
\[
  \operatorname{weightedAncestorCoeffSum}(G,R,Q).
\]

The second term \(U_{2,i}\) collects the complementary interactions, where the
source cell is below the indicator cell.  Its coefficients have the schematic
form
\begin{equation}\label{eq:u2-coeff}
  (U_{2,i})_J
  =
   \Theta_i c_J\,\Tcal_i^{<}(J),
\end{equation}
where \(\Tcal_i^{<}(J)\) is the strict ancestor tower of the indicator
representation above \(J\).  More explicitly, if \(J\in\Pcal^j\), then
\begin{equation}\label{eq:strict-indicator-tower}
  \Tcal_i^{<}(J)
  =
  \sum_{\ell<j}
  \sum_{\substack{Q\in\Pcal^\ell\\ J\subset Q}}
  b^i_Q\,|Q|^{s-\frac1p}.
\end{equation}
The word ``strict'' means that the level \(\ell=j\) is excluded. Since this representation is positive we have 

$$\| \Tcal_i^{<}(J)\|\leq \|1_{\Omega_i}\|_\infty\leq 1.$$


 In the Lean
formalization this is
\[
  \operatorname{strictWeightedAncestorCoeffSum}(G,R_i,J),
\]
where \(R_i\) is the chosen indicator representation of \(\1_{\Om_i}\).

Define 

$$U_1 = \sum_{i\in \Lambda}  U_{1,i}$$
and
$$U_2 = \sum_{i\in \Lambda}  U_{2,i}$$

We have 
 $$(U_{1})_Q
  =
  \sum_{i\in \Lambda}  \Theta_i b^i_Q\,\Tcal_R(Q),$$
$$(U_{2})_J 
  =
  \sum_{i\in \Lambda}   \Theta_i c_J\,\Tcal_i^{<}(J),$$
  
  
\subsection*{The \(U_1\) estimate}

The cost at level $k$ of $U_1$ is, by Minkovsky

\begin{align} \left(  \sum_{Q\in \mathcal{P}^k} \|\sum_{i\in \Lambda} \Theta_i b^i_Q\,\Tcal_R(Q)\|^p \right)^{1/p} \\
\leq \left(  \sum_{Q\in \mathcal{P}^k} \|\sum_{ i\in \Lambda } \Theta_i b^i_Q\,\Tcal_R(Q)\|^p \right)^{1/p} \\
\leq  M \left(  \sum_{Q\in \mathcal{P}^k} \|\sum_{ i\in \Lambda}  \Theta_i b^i_Q\|^p \right)^{1/p}\\
\leq  M \sum_{ i\in \Lambda}  \| \Theta_i \|\left(  \sum_{Q\in \mathcal{P}^k} \|b^i_Q\|^p \right)^{1/p}
\end{align} 

so 


\begin{align}\left(  \sum_k \left(  \sum_{Q\in \mathcal{P}^k} \|\sum_{i\in \Lambda}b^i_Q\,\Tcal_R(Q)\|^p \right)^{q/p}\right)^{1/q}  \\
\leq  M \left( \sum_k \left( \sum_{\i\in \Lambda} \|\Theta_i \|  \left(  \sum_{Q\in \mathcal{P}^k} \|b^i_Q\|^p \right)^{1/p} \right)^{q} \right)^{1/q} \\
\leq  M \sum_{\i\in \Lambda}  \|\Theta_i \| \left( \sum_k \left(    \sum_{Q\in \mathcal{P}^k} \|b^i_Q\|^p \right)^{q/p}  \right)^{1/q} \\
\leq M \sum_{i\in \Lambda} \|\Theta_i \| \Ccal_i.
\end{align} 

For a fixed \(Q\in\Pcal^j\), suppose
\[
  \Th_i (U_{1,i})_Q\neq0.
\]
Then \(b^i_Q\neq0\), hence \(Q\subset\Om_i\).  Also, if
\(\Tcal_R(Q)\neq0\), then some ancestor \(P\supset Q\) has \(c_P\neq0\).
Thus \(P\cap\Om_i\neq\emptyset\).  Applying the overlap hypothesis
\eqref{eq:overlap} to this active cell \(P\) gives
\[
  \abs{\Th_i}(1+\Ccal_i)\le N.
\]
Since \(\abs{b^i_Q}\le \Ccal_i\le1+\Ccal_i\), we obtain the local bound
\begin{equation}\label{eq:u1-local}
  \abs{\Th_i}\abs{b^i_Q}
  \le N.
\end{equation}
Using \eqref{eq:u1-coeff},
\[
  \abs{\Th_i(U_{1,i})_Q}
  =
  \abs{\Th_i}\abs{b^i_Q}\abs{\Tcal_R(Q)}
  \le
  N\abs{\Tcal_R(Q)}.
\]

Now use disjointness of the domains.  For a fixed \(Q\), at most one index
\(i\) can have \(b^i_Q\neq0\), because \(b^i_Q\neq0\) implies
\(Q\subset\Om_i\), and the \(\Om_i\)'s are pairwise disjoint.  Therefore, for
every finite \(\Gamma\subset\Lambda\),
\begin{align}
  \sum_{i\in\Gamma}\abs{\Th_i}^p
  \sum_{Q\in\Pcal^j}\abs{(U_{1,i})_Q}^p
  &=
  \sum_{Q\in\Pcal^j}
  \sum_{i\in\Gamma}
  \abs{\Th_i(U_{1,i})_Q}^p
  \nonumber\\
  &\le
  \sum_{Q\in\Pcal^j}
  N^p\abs{\Tcal_R(Q)}^p
  \nonumber\\
  &=
  N^p T_j(R)^p.
  \label{eq:u1-level}
\end{align}

\subsection*{The \(U_2\) estimate}

The \(U_2\) term is controlled by the source coefficients themselves.  Indeed,
the strict indicator tower \(\Tcal_i^{<}(J)\) is nonzero only when an ancestor
of \(J\) belongs to the canonical decomposition of \(\Om_i\).  Thus, if
\(\Th_i(U_{2,i})_J\neq0\), then some source cell meeting the corresponding
indicator ancestor also meets \(\Om_i\), and the overlap hypothesis again gives
\[
  \abs{\Th_i}(1+\Ccal_i)\le N.
\]
The canonical indicator construction gives
\[
  \abs{\Tcal_i^{<}(J)}\le 1+\Ccal_i.
\]
Hence
\[
  \abs{\Th_i(U_{2,i})_J}
  \le
  N\abs{c_J}.
\]
As before, disjointness implies that for each \(J\) at most one \(i\) may
contribute.  Consequently,
\begin{align}
  \sum_{i\in\Gamma}\abs{\Th_i}^p
  \sum_{J\in\Pcal^j}\abs{(U_{2,i})_J}^p
  &\le
  N^p
  \sum_{J\in\Pcal^j}\abs{c_J}^p
  \nonumber\\
  &=
  N^p A_j(R)^p.
  \label{eq:u2-level}
\end{align}

\subsection*{Combining \(U_1\) and \(U_2\)}

The product representation for the finite partial sum
\[
  h_\Gamma
  =
  \sum_{i\in\Gamma}\Th_i\1_{\Om_i}f
\]
has level \(j\) coefficient power bounded by the weighted sum of
\(U_{1,i}+U_{2,i}\).  Using
\[
  \abs{x+y}^p\le 2^{p-1}\bigl(\abs{x}^p+\abs{y}^p\bigr),
\]
together with \eqref{eq:u1-level} and \eqref{eq:u2-level}, we get
\begin{align}
  \sum_{i\in\Gamma}\abs{\Th_i}^p
  \sum_{Q\in\Pcal^j}
  \abs{(U_{1,i}+U_{2,i})_Q}^p
  &\le
  2^{p-1}N^p
  \bigl(T_j(R)^p+A_j(R)^p\bigr).
  \label{eq:level-final}
\end{align}
This is the non-uniform level-by-level estimate.  Notice that the factor
\(\abs{\Th_i}\Ccal_i\) is used locally before summing in \(Q\); it is not
replaced by a global counting estimate.

\section*{The Tower Estimate}

The remaining analytic input is the standard tower estimate for canonical
Souza representations:
\begin{equation}\label{eq:tower-estimate}
  \begin{cases}
  \displaystyle
  \left(
    \sum_{j\ge0} T_j(R)^q
  \right)^{1/q}
  \le
  C_{\mathrm{tow}}
  \bigl(\operatorname{pqCost}(R)+M\bigr),
  & q<\infty,\\[2ex]
  \displaystyle
  \sup_{j\ge0}T_j(R)
  \le
  C_{\mathrm{tow}}
  \bigl(\operatorname{pqCost}(R)+M\bigr),
  & q=\infty.
  \end{cases}
\end{equation}
This is the paper-level form of the ``tower term''.  It is a discrete Hardy
estimate along the ancestor chains of the grid.  The point of
\eqref{eq:tower-estimate} is that one controls the whole sequence
\((T_j(R))_j\), not merely the pointwise bound
\(\abs{\Tcal_R(Q)}\le M\).  The latter would introduce a false factor counting
the number of cells on level \(j\).

\section*{Mixed Cost Estimate}

Let \(S_\Gamma\) be the representation of \(h_\Gamma\) obtained from the above
product blocks.  From \eqref{eq:level-final},
\[
  \operatorname{levelCoeffPower}_j(S_\Gamma)
  \le
  2^{p-1}N^p
  \bigl(T_j(R)^p+A_j(R)^p\bigr).
\]
Taking \(p\)-roots gives
\[
  \operatorname{levelCoeffPower}_j(S_\Gamma)^{1/p}
  \le
  2^{1-\frac1p}N
  \bigl(T_j(R)^p+A_j(R)^p\bigr)^{1/p}.
\]
Since
\[
  \bigl(a^p+b^p\bigr)^{1/p}\le a+b
  \qquad (a,b\ge0),
\]
we obtain
\[
  \operatorname{levelCoeffPower}_j(S_\Gamma)^{1/p}
  \le
  2^{1-\frac1p}N
  \bigl(T_j(R)+A_j(R)\bigr).
\]
Taking the \(\ell^q\) norm in \(j\), using Minkowski's inequality and
\eqref{eq:tower-estimate},
\begin{align}
  \operatorname{pqCost}(S_\Gamma)
  &\le
  2^{1-\frac1p}N
  \left(
    \norm{(T_j(R))_j}_{\ell^q}
    +
    \norm{(A_j(R))_j}_{\ell^q}
  \right)
  \nonumber\\
  &\le
  C_1 N
  \bigl(\operatorname{pqCost}(R)+M\bigr),
  \label{eq:finite-cost}
\end{align}
where \(C_1\) depends only on \(s,p,q\) and the grid constants.

\section*{Pointwise and \(L^\infty\) Estimates}

For a fixed point \(x\), pairwise disjointness of the domains implies that at
most one indicator \(\1_{\Om_i}(x)\) is nonzero.  If \(f(x)\neq0\), choose an
active source cell \(P\) containing \(x\).  Then every \(i\) with
\(\1_{\Om_i}(x)\neq0\) satisfies \(P\cap\Om_i\neq\emptyset\).  Hence
\eqref{eq:overlap} gives
\[
  \sum_{i\in\Lambda}
  \abs{\Th_i\1_{\Om_i}(x)}
  \le N.
\]
Therefore
\[
  \abs{h(x)}
  =
  \left|
    \sum_{i\in\Lambda}
    \Th_i\1_{\Om_i}(x)f(x)
  \right|
  \le
  N\abs{f(x)}
  \le
  NM
\]
for a.e. \(x\).  Thus
\begin{equation}\label{eq:linfty}
  \norm{h}_{L^\infty}\le NM.
\end{equation}

\section*{Passage to the Infinite Sum}

Let \((\Gamma_n)_n\) be an increasing sequence of finite subsets of
\(\Lambda\) exhausting \(\Lambda\), and set
\[
  h_n=\sum_{i\in\Gamma_n}\Th_i\1_{\Om_i}f .
\]
The estimates \eqref{eq:finite-cost} and \eqref{eq:linfty} are uniform in
\(n\).  By the compactness/limit theorem for bounded families of Besov
representations, there is a subsequence of the representations \(S_{\Gamma_n}\)
converging to a representation \(S\) of
\[
  h=\sum_{i\in\Lambda}\Th_i\1_{\Om_i}f
\]
with
\[
  \operatorname{pqCost}(S)
  \le
  C_1N\bigl(\operatorname{pqCost}(R)+M\bigr).
\]
The support condition is closed under this limit: if a coefficient of \(S\) on
a cell \(Q\) is nonzero, then \(Q\subset\Om_i\) for some \(i\in\Lambda\).
Combining this with \eqref{eq:linfty}, and increasing the constant if
necessary, gives
\[
  \operatorname{pqCost}(S)+\norm{h}_{L^\infty}
  \le
  C_{\mathrm{na}}N
  \bigl(\operatorname{pqCost}(R)+M\bigr).
\]
This proves \eqref{eq:main-bound}.

\begin{remark}
In the Lean development, the already formalized level estimate is exactly
\eqref{eq:level-final}.  The term
\[
  T_j(R)^p
  =
  \sum_{Q\in\Pcal^j}
  \left|
    \operatorname{weightedAncestorCoeffSum}(G,R,Q)
  \right|^p
\]
is the tower term.  The remaining formal analytic step is the mixed
\(\ell^q\) estimate \eqref{eq:tower-estimate}; replacing \(T_j(R)\) by the
pointwise bound \(M\) times the number of level cells is not the desired
argument.
\end{remark}

\end{document}
